% ==============================================================================
% Relatório Técnico - Academic Events API
% Curso de Desenvolvimento em C# - Módulo 3
% Grupo 6: Thailsson, Stevão, Márcio, Allef, Juliana
% ==============================================================================

\documentclass[11pt, a4paper, oneside]{article}

% ==============================================================================
% Encoding e idioma
% ==============================================================================
\usepackage[T1]{fontenc}
\usepackage[utf8]{inputenc}
\usepackage[brazil]{babel}

% ==============================================================================
% Layout
% ==============================================================================
\usepackage{geometry}
\usepackage{lmodern}
\usepackage{microtype}
\usepackage{setspace}
\usepackage{parskip}
\usepackage{fancyhdr}
\usepackage{lastpage}
\usepackage{titlesec}
\usepackage{tocloft}
\usepackage{float}
\usepackage{xurl}
\usepackage{amsmath}

% ==============================================================================
% Gráficos, cores e diagramas
% ==============================================================================
\usepackage{graphicx}
\usepackage{xcolor}
\usepackage{tikz}
\usetikzlibrary{positioning, arrows.meta}

% ==============================================================================
% Tabelas e listas
% ==============================================================================
\usepackage{booktabs}
\usepackage{longtable}
\usepackage{tabularx}
\usepackage{array}
\usepackage{multirow}
\usepackage{colortbl}
\usepackage{ragged2e}
\usepackage{enumitem}

% ==============================================================================
% Caixas e código
% ==============================================================================
\usepackage{tcolorbox}
\tcbuselibrary{skins,breakable}
\usepackage{listings}
\usepackage{hyperref}

% ==============================================================================
% Geometria da página
% ==============================================================================
\geometry{
  a4paper,
  top=25mm, bottom=25mm,
  left=25mm, right=20mm,
  headheight=15mm, headsep=6mm,
  footskip=12mm
}

\setlength{\emergencystretch}{3em}
\setlength{\hfuzz}{2pt}
\setlength{\LTleft}{0pt}
\setlength{\LTright}{0pt}
\setlength{\tabcolsep}{4pt}
\Urlmuskip=0mu plus 2mu\relax

% ==============================================================================
% Cores (identidade do projeto)
% ==============================================================================
\definecolor{apiblue}{RGB}{30, 58, 138}
\definecolor{apilightblue}{RGB}{59, 130, 246}
\definecolor{apidark}{RGB}{15, 23, 42}
\definecolor{apilight}{RGB}{239, 246, 255}
\definecolor{tablehead}{RGB}{30, 58, 138}
\definecolor{tablealt}{RGB}{239, 246, 255}
\definecolor{codebg}{RGB}{248, 250, 252}
\definecolor{alertcolor}{RGB}{180, 100, 0}
\definecolor{successcolor}{RGB}{22, 101, 52}
\definecolor{successlight}{RGB}{220, 252, 231}

% ==============================================================================
% Hyperref
% ==============================================================================
\hypersetup{
  colorlinks  = true,
  linkcolor   = apiblue,
  urlcolor    = apilightblue,
  citecolor   = apidark,
  pdftitle    = {Relatório Técnico - Academic Events API - Grupo 6},
  pdfauthor   = {Thailsson Clementino; Stevão Whinter; Márcio Franklin; Allef Oliveira; Juliana Ballin},
  pdfsubject  = {API REST em ASP.NET Core para gerenciamento de eventos acadêmicos},
  pdfkeywords = {CSharp, ASP.NET Core, REST API, JWT, EF Core, PostgreSQL, DDD}
}

% ==============================================================================
% Tipos de coluna
% ==============================================================================
\newcolumntype{J}[1]{>{\justifying\arraybackslash\small}p{#1}}
\newcolumntype{L}[1]{>{\raggedright\arraybackslash\small}p{#1}}
\newcolumntype{C}[1]{>{\centering\arraybackslash\small}p{#1}}
\newcolumntype{Y}{>{\justifying\arraybackslash\small}X}
\renewcommand{\arraystretch}{1.3}

% ==============================================================================
% Formatação de seções
% ==============================================================================
\titleformat{\section}[block]
  {\Large\bfseries\color{apiblue}}
  {\thesection}{1em}{}
  [\vspace{1pt}{\color{apilightblue}\hrule height 1.5pt}\vspace{4pt}]

\titleformat{\subsection}[block]
  {\large\bfseries\color{apiblue}}
  {\thesubsection}{1em}{}

\titleformat{\subsubsection}[block]
  {\normalsize\bfseries\color{apidark}}
  {\thesubsubsection}{1em}{}

\titlespacing{\section}{0pt}{18pt}{6pt}
\titlespacing{\subsection}{0pt}{12pt}{4pt}
\titlespacing{\subsubsection}{0pt}{8pt}{2pt}

% ==============================================================================
% Configuração de código
% ==============================================================================
\lstset{
  backgroundcolor=\color{codebg},
  basicstyle=\ttfamily\footnotesize,
  breaklines=true,
  frame=single,
  rulecolor=\color{apilightblue},
  numbers=none,
  keywordstyle=\color{apiblue}\bfseries,
  stringstyle=\color{successcolor},
  commentstyle=\color{gray}\itshape,
  showstringspaces=false,
  extendedchars=true,
  inputencoding=utf8,
  literate=
    {á}{{\'a}}1 {à}{{\`a}}1 {â}{{\^a}}1 {ã}{{\~a}}1
    {Á}{{\'A}}1 {À}{{\`A}}1 {Â}{{\^A}}1 {Ã}{{\~A}}1
    {é}{{\'e}}1 {ê}{{\^e}}1 {É}{{\'E}}1 {Ê}{{\^E}}1
    {í}{{\'i}}1 {Í}{{\'I}}1
    {ó}{{\'o}}1 {ô}{{\^o}}1 {õ}{{\~o}}1
    {Ó}{{\'O}}1 {Ô}{{\^O}}1 {Õ}{{\~O}}1
    {ú}{{\'u}}1 {Ú}{{\'U}}1
    {ç}{{\c{c}}}1 {Ç}{{\c{C}}}1
}

% ==============================================================================
% Cabeçalho e rodapé
% ==============================================================================
\pagestyle{fancy}
\fancyhf{}
\fancyhead[L]{\small\color{apiblue}\textbf{Academic Events API}}
\fancyhead[R]{\small\color{apidark}Grupo 6 - Módulo Desenvolvimento em C\#}
\fancyfoot[C]{\small\color{apidark}Página \thepage\ de \pageref{LastPage}}
\renewcommand{\headrulewidth}{1.5pt}
\renewcommand{\headrule}{\hbox to\headwidth{\color{apilightblue}\leaders\hrule height \headrulewidth\hfill}}

% ==============================================================================
% Caixas personalizadas
% ==============================================================================
\tcbset{
  requisito/.style={
    enhanced, breakable,
    colback=apilight, colframe=apiblue,
    fonttitle=\bfseries\color{white},
    attach boxed title to top left={yshift=-2mm, xshift=4mm},
    boxed title style={colback=apiblue, rounded corners},
    sharp corners=south
  },
  destaque/.style={
    enhanced,
    colback=successlight, colframe=successcolor,
    fonttitle=\bfseries\color{successcolor},
    sharp corners
  }
}

% ==============================================================================
% Início do documento
% ==============================================================================
\begin{document}

% ==============================================================================
% CAPA
% ==============================================================================
\begin{titlepage}
\centering
\vspace*{20mm}

{\color{apiblue}\rule{\linewidth}{2pt}}\\[6mm]

{\Huge\bfseries\color{apiblue} Academic Events API}\\[4mm]
{\Large\color{apidark} Relatório Técnico - Trabalho Final}\\[3mm]
{\large\color{apidark} Módulo 3: Desenvolvimento em C\texttt{\#}}\\[8mm]

{\color{apilightblue}\rule{\linewidth}{1pt}}\\[10mm]

\begin{tabular}{ll}
\textbf{Disciplina:} & Desenvolvimento em C\# - Módulo 3 \\[2mm]
\textbf{Grupo:}      & Grupo 6 \\[2mm]
\textbf{Data:}       & Junho de 2026 \\[2mm]
\end{tabular}

\vspace{10mm}

\begin{tcolorbox}[requisito, title={Integrantes do Grupo 6}, width=0.85\linewidth]
\begin{tabular}{ll}
1. & Thailsson Clementino de Andrade \\
2. & Stevão Whinter Marques de Andrade \\
3. & Márcio Franklin de Oliveira Lima \\
4. & Allef Oliveira Ramos \\
5. & Juliana Ballin Lima \\
\end{tabular}
\end{tcolorbox}

\vfill

{\color{apiblue}\rule{\linewidth}{2pt}}\\[4mm]
{\small\color{apidark} Repositório: \url{https://github.com/JulianaBallin/Academic-Events-API}}
\end{titlepage}

% ==============================================================================
% SUMÁRIO
% ==============================================================================
\tableofcontents
\newpage

% ==============================================================================
\section{Introdução}
% ==============================================================================

O presente relatório descreve o desenvolvimento da \textbf{Academic Events API}, uma API REST criada como trabalho final do Módulo 3 do curso de Desenvolvimento em C\#. O projeto implementa um sistema de gerenciamento de eventos acadêmicos, como palestras, workshops e seminários, seguindo as boas práticas de desenvolvimento com ASP.NET Core 8.

O objetivo principal foi aplicar os conceitos ensinados ao longo do módulo: orientação a objetos, DTOs, services, repositories, acesso a banco de dados, autenticação e autorização, e organização profissional de código em camadas separadas.

A plataforma permite que usuários se cadastrem, criem e gerenciem eventos acadêmicos, realizem inscrições, publiquem comentários e expressem reações a conteúdos de interesse, funcionando de forma análoga a uma rede social voltada para o ambiente acadêmico.

% ==============================================================================
\section{Problema e Solução}
% ==============================================================================

\subsection{O Problema}

Eventos acadêmicos, como palestras, workshops, minicursos e seminários, são parte essencial da vida universitária. No entanto, a organização e divulgação desses eventos muitas vezes acontece de forma fragmentada: grupos de WhatsApp, cartazes físicos e e-mails avulsos dificultam o acompanhamento das oportunidades disponíveis e o controle de inscrições.

\subsection{Nossa Solução}

A Academic Events API centraliza o ciclo de vida completo de um evento acadêmico em uma única plataforma, expondo endpoints REST documentados via Swagger. Qualquer aplicação cliente, seja web, mobile ou desktop, pode consumir a API para:

\begin{itemize}[noitemsep]
  \item Listar e filtrar eventos por status ou organizador
  \item Criar e gerenciar eventos (organizadores autenticados)
  \item Realizar e acompanhar inscrições
  \item Comentar e reagir a eventos de interesse
\end{itemize}

% ==============================================================================
\section{Tecnologias Utilizadas}
% ==============================================================================

\begin{longtable}{L{4.5cm} L{4cm} L{6cm}}
\toprule
\rowcolor{tablehead}
\color{white}\textbf{Tecnologia} & \color{white}\textbf{Versão} & \color{white}\textbf{Finalidade} \\
\midrule
\endhead
ASP.NET Core & 8.0 & Framework principal da API REST \\
\rowcolor{tablealt}
Entity Framework Core & 8.0.11 & ORM para acesso ao banco de dados \\
Npgsql (PostgreSQL) & 8.0.11 & Driver de conexão com PostgreSQL \\
\rowcolor{tablealt}
JWT Bearer Token & 8.0.15 & Autenticação e autorização via token \\
BCrypt.Net-Next & 4.2.0 & Hash seguro de senhas dos usuários \\
\rowcolor{tablealt}
Swashbuckle / Swagger & 6.6.2 & Documentação automática dos endpoints \\
PostgreSQL & 15 & Banco de dados relacional (Docker) \\
\rowcolor{tablealt}
Docker Compose & - & Ambiente de banco de dados isolado \\
\bottomrule
\end{longtable}

% ==============================================================================
\section{Arquitetura do Sistema}
% ==============================================================================

\subsection{Visão Geral - Clean Architecture em Camadas}

O projeto é organizado em cinco projetos C\# separados dentro de uma única solution, seguindo os princípios de separação de responsabilidades e inversão de dependência:

\begin{tcolorbox}[requisito, title={Estrutura da Solution}]
\begin{itemize}[noitemsep]
  \item \textbf{AcademicEvents.Domain} - Entidades, enums, interfaces de repository (sem dependências externas)
  \item \textbf{AcademicEvents.Exceptions} - Exceções customizadas com semântica de negócio
  \item \textbf{AcademicEvents.Application} - DTOs, interfaces de service, implementações de service
  \item \textbf{AcademicEvents.Infrastructure} - DbContext, EF Core, implementações de repository
  \item \textbf{AcademicEvents.API} - Controllers, JWT, Swagger, Program.cs (ponto de entrada)
\end{itemize}
\end{tcolorbox}

\subsection{Fluxo de Dependências entre Projetos}

O gráfico abaixo ilustra como os projetos se relacionam, respeitando a regra de que as camadas internas nunca dependem das externas:

\begin{center}
\begin{tikzpicture}[
  node distance=1.4cm,
  every node/.style={
    rectangle, rounded corners=4pt,
    minimum width=4.5cm, minimum height=0.9cm,
    font=\small\bfseries, align=center
  }
]
  \node[fill=apiblue!20, draw=apiblue] (api) {AcademicEvents.API};
  \node[fill=apiblue!10, draw=apiblue, below=of api] (app) {AcademicEvents.Application};
  \node[fill=apiblue!10, draw=apiblue, right=2cm of app] (infra) {AcademicEvents.Infrastructure};
  \node[fill=gray!15, draw=gray, below=of app] (domain) {AcademicEvents.Domain};
  \node[fill=gray!15, draw=gray, right=2cm of domain] (exc) {AcademicEvents.Exceptions};

  \draw[->, thick, color=apiblue] (api) -- (app);
  \draw[->, thick, color=apiblue] (api) -- (infra);
  \draw[->, thick, color=apiblue] (app) -- (domain);
  \draw[->, thick, color=apiblue] (infra) -- (domain);
  \draw[->, thick, color=gray] (app) -- (exc);
  \draw[->, thick, color=gray] (api) -- (exc);
\end{tikzpicture}
\end{center}

\subsection{Por que essa Arquitetura?}

\begin{itemize}
  \item \textbf{Domain sem dependências:} as entidades e interfaces ficam isoladas, podendo ser testadas sem infraestrutura
  \item \textbf{Application só depende de Domain:} os services não sabem como os dados são persistidos. Apenas chamam interfaces
  \item \textbf{Infrastructure implementa interfaces do Domain:} a inversão de dependência permite trocar o banco (ex: SQLite para testes) sem mudar a lógica de negócio
  \item \textbf{API orquestra tudo via DI:} o Program.cs conecta as implementações concretas usando extension methods de injeção de dependência
\end{itemize}

% ==============================================================================
\section{Modelagem do Domínio}
% ==============================================================================

\subsection{Entidades}

O domínio é composto por cinco entidades principais com os respectivos relacionamentos:

\begin{longtable}{L{3cm} L{5cm} L{6.5cm}}
\toprule
\rowcolor{tablehead}
\color{white}\textbf{Entidade} & \color{white}\textbf{Campos principais} & \color{white}\textbf{Descrição} \\
\midrule
\endhead
User & Id, Nome, Email, SenhaHash, CriadoEm & Usuário da plataforma. Pode organizar eventos, se inscrever, comentar e reagir \\
\rowcolor{tablealt}
Event & Id, Titulo, Descricao, DataInicio, DataFim, Local, Status, OrganizadorId & Evento acadêmico criado por um organizador. Ciclo: Rascunho, Publicado, Cancelado, Concluído \\
Registration & Id, UsuarioId, EventoId, Status & Inscrição de um usuário em um evento. Índice único impede duplicatas \\
\rowcolor{tablealt}
Comment & Id, UsuarioId, EventoId, Conteudo & Comentário do usuário em um evento \\
Reaction & Id, UsuarioId, EventoId, Tipo & Reação (curtida) de um usuário a um evento. Um usuário só pode ter uma reação por evento \\
\bottomrule
\end{longtable}

\subsection{Enumeradores}

\begin{longtable}{L{4cm} L{10.5cm}}
\toprule
\rowcolor{tablehead}
\color{white}\textbf{Enum} & \color{white}\textbf{Valores} \\
\midrule
\endhead
StatusEvento & Rascunho, Publicado, Cancelado, Concluído \\
\rowcolor{tablealt}
StatusInscricao & Pendente, Confirmada, Cancelada \\
TipoReacao & Curtir, Adorei, Interessante, VouParticipar \\
\bottomrule
\end{longtable}

% ==============================================================================
\section{Autenticação e Autorização}
% ==============================================================================

\subsection{JWT Bearer Token}

A autenticação utiliza o padrão \textbf{JWT (JSON Web Token)} com algoritmo HMAC-SHA256. O fluxo é:

\begin{enumerate}
  \item Usuário envia email e senha para \texttt{POST /api/auth/login}
  \item O \texttt{AuthService} valida as credenciais com BCrypt
  \item Um token JWT é gerado com as \textit{claims}: \texttt{NameIdentifier} (id), \texttt{Email} e \texttt{Name}
  \item O token é retornado e deve ser enviado no header \texttt{Authorization: Bearer \{token\}}
  \item O middleware do ASP.NET valida o token automaticamente em todas as rotas com \texttt{[Authorize]}
\end{enumerate}

\subsection{Configuração do Token}

\begin{lstlisting}[language={[Sharp]C}]
// appsettings.json
"Jwt": {
  "Key": "chave-secreta-mínimo-32-caracteres-aqui",
  "Issuer": "AcademicEventsAPI",
  "Audience": "AcademicEventsClientes",
  "ExpiresInHours": 8
}
\end{lstlisting}

\subsection{Hash de Senha}

As senhas dos usuários nunca são armazenadas em texto puro. O \textbf{BCrypt} é usado para gerar um hash com salt automático, tornando ataques de dicionário impraticáveis mesmo que o banco seja comprometido.

% ==============================================================================
\section{Endpoints da API}
% ==============================================================================

\subsection{Autenticação (público)}

\begin{longtable}{L{1.5cm} L{4cm} L{3.5cm} L{5.5cm}}
\toprule
\rowcolor{tablehead}
\color{white}\textbf{Método} & \color{white}\textbf{Rota} & \color{white}\textbf{Retorno} & \color{white}\textbf{Descrição} \\
\midrule
\endhead
POST & /api/auth/register & 200, 400 & Cadastra usuário e retorna JWT \\
\rowcolor{tablealt}
POST & /api/auth/login & 200, 401 & Autentica usuário e retorna JWT \\
GET & /api/me & 200, 401 & Dados do usuário autenticado (JWT) \\
\bottomrule
\end{longtable}

\subsection{Eventos}

\begin{longtable}{L{1.5cm} L{5cm} L{2.5cm} L{5.5cm}}
\toprule
\rowcolor{tablehead}
\color{white}\textbf{Método} & \color{white}\textbf{Rota} & \color{white}\textbf{Auth} & \color{white}\textbf{Descrição} \\
\midrule
\endhead
GET & /api/events & Não & Lista todos os eventos \\
\rowcolor{tablealt}
GET & /api/events?status=X & Não & Filtra por status \\
GET & /api/events/\{id\} & Não & Busca evento por id \\
\rowcolor{tablealt}
GET & /api/events?organizadorId=X & Não & Filtra por organizador \\
GET & /api/events?status=X\&organizadorId=Y & Não & Combina filtros de status e organizador \\
\rowcolor{tablealt}
GET & /api/events/meus & Sim & Eventos do organizador autenticado \\
POST & /api/events & Sim & Cria novo evento \\
\rowcolor{tablealt}
PUT & /api/events/\{id\} & Sim & Atualiza evento (só o organizador) \\
DELETE & /api/events/\{id\} & Sim & Remove evento (só o organizador) \\
\bottomrule
\end{longtable}

\subsection{Inscrições, Comentários e Reações}

\begin{longtable}{L{1.5cm} L{5.5cm} L{2cm} L{5cm}}
\toprule
\rowcolor{tablehead}
\color{white}\textbf{Método} & \color{white}\textbf{Rota} & \color{white}\textbf{Auth} & \color{white}\textbf{Descrição} \\
\midrule
\endhead
POST & /api/registrations & Sim & Inscreve usuário em evento \\
\rowcolor{tablealt}
GET & /api/registrations/me & Sim & Minhas inscrições \\
DELETE & /api/registrations/\{id\} & Sim & Cancela inscrição \\
\midrule
GET & /api/comments?eventoId=X & Não & Comentários de um evento \\
\rowcolor{tablealt}
POST & /api/comments & Sim & Adiciona comentário \\
DELETE & /api/comments/\{id\} & Sim & Remove comentário (só o autor) \\
\midrule
GET & /api/reactions?eventoId=X & Não & Reações de um evento \\
\rowcolor{tablealt}
POST & /api/reactions & Sim & Adiciona reação \\
DELETE & /api/reactions/\{id\} & Sim & Remove reação (só o autor) \\
\bottomrule
\end{longtable}

% ==============================================================================
\section{Camada de Infrastructure}
% ==============================================================================

\subsection{DbContext e Configurações EF Core}

O \texttt{AcademicEventsDbContext} configura os mapeamentos via \textit{Fluent API} no método \texttt{OnModelCreating}:

\begin{itemize}[noitemsep]
  \item \textbf{Índice único em User.Email}: impede dois usuários com o mesmo email
  \item \textbf{Índice único em (UsuarioId, EventoId) na Registration}: impede inscrição duplicada
  \item \textbf{Índice único em (UsuarioId, EventoId) na Reaction}: impede reação duplicada
  \item \textbf{DeleteBehavior.Restrict} em Event para User: impede deletar usuário que tem eventos
  \item \textbf{DeleteBehavior.Cascade} nos demais: ao remover usuário ou evento, registros dependentes são apagados
\end{itemize}

\subsection{Repository Pattern}

Cada entidade tem seu próprio repository que abstrai o acesso ao banco:

\begin{lstlisting}[language={[Sharp]C}]
// exemplo de consulta com Include para evitar N+1 queries
public async Task<List<Event>> GetAllAsync()
{
    return await _context.Events
        .Include(e => e.Organizador)
        .OrderByDescending(e => e.DataInicio)
        .ToListAsync();
}
\end{lstlisting}

O uso de \texttt{Include} garante que o nome do organizador seja carregado junto com o evento em uma única query SQL, evitando o problema N+1.

\subsection{Banco de Dados via Docker}

O ambiente de banco de dados é provisionado via \texttt{docker-compose.yml}:

\begin{lstlisting}
docker compose up -d   # sobe o PostgreSQL 15 na porta 5432
\end{lstlisting}

A inicialização do schema usa \texttt{EnsureCreated()} no startup da API, criando as tabelas automaticamente se não existirem.

% ==============================================================================
\section{Regras de Negócio Implementadas}
% ==============================================================================

\begin{longtable}{L{5cm} L{9.5cm}}
\toprule
\rowcolor{tablehead}
\color{white}\textbf{Regra} & \color{white}\textbf{Implementação} \\
\midrule
\endhead
Email único & Verificado no AuthService antes de inserir no banco (400 BadRequest) \\
\rowcolor{tablealt}
Senha segura (min 6 chars) & Validação de DataAnnotations no RegisterRequest \\
Credenciais inválidas & Mesma mensagem independente de email ou senha errados (segurança) \\
\rowcolor{tablealt}
DataFim maior que DataInicio & Validado no EventService ao criar e atualizar evento (400 BadRequest) \\
Só organizador edita evento & Verificação de OrganizadorId no EventService (403 Forbidden) \\
\rowcolor{tablealt}
Inscrição duplicada & Verificação no RegistrationService antes do banco (400 BadRequest) \\
Evento existente em inscrição & RegistrationService retorna 404 quando o evento informado não existe \\
\rowcolor{tablealt}
Reação única por evento & Verificação no ReactionService antes do banco (400 BadRequest) \\
Evento existente em comentários e reações & CommentService e ReactionService retornam 404 quando o evento informado não existe \\
\rowcolor{tablealt}
Só autor deleta comentário & Verificação de UsuarioId no CommentService (403 Forbidden) \\
Respostas HTTP corretas & 200/201 sucesso, 400 dados inválidos, 401 sem auth, 403 sem permissão, 404 não encontrado \\
\bottomrule
\end{longtable}

% ==============================================================================
\section{Dificuldades e Decisões Técnicas}
% ==============================================================================

\subsection{Versões dos Pacotes NuGet}

A combinação de versões dos pacotes do Entity Framework Core para .NET 8 exigiu atenção especial. A versão \texttt{8.0.11} do \texttt{Npgsql.EntityFrameworkCore.PostgreSQL} foi escolhida por ser a mais estável e compatível com o .NET 8 e o \texttt{Microsoft.EntityFrameworkCore.Design 8.0.11}.

\subsection{Forbid() vs StatusCode(403)}

O método \texttt{Forbid()} do \texttt{ControllerBase} do ASP.NET Core não aceita uma mensagem de texto. Ele espera um nome de esquema de autenticação. Para retornar 403 com mensagem de erro no corpo da resposta, o correto é usar \texttt{StatusCode(403, mensagem)}.

\subsection{Application não depende de Infrastructure}

A camada Application só referencia o Domain (onde estão as interfaces de repository). A Infrastructure é registrada no DI pela API via extension method, resolvendo as implementações concretas em tempo de execução. Isso garante a inversão de dependência correta.

\subsection{Validação em Dois Níveis}

Para inscrições e reações duplicadas, a validação acontece tanto no service (verificação prévia antes do banco) quanto no banco (índice único). O service lança uma exceção customizada com mensagem amigável; o banco é a última linha de defesa contra condições de corrida.

% ==============================================================================
\section{Como Executar o Projeto}
% ==============================================================================

\begin{tcolorbox}[destaque, title={Requisitos}]
\begin{itemize}[noitemsep]
  \item .NET 8 SDK
  \item Docker e Docker Compose
  \item (Opcional) dotnet-ef para migrations: \texttt{dotnet tool install -{}-global dotnet-ef}
\end{itemize}
\end{tcolorbox}

\begin{lstlisting}[language=bash]
# 1. clonar o repositório
git clone https://github.com/JulianaBallin/Academic-Events-API.git
cd Academic-Events-API

# 2. subir o banco PostgreSQL via Docker
docker compose up -d

# 3. iniciar a API (o EnsureCreated cria as tabelas automaticamente)
cd AcademicEvents.API
dotnet run
\end{lstlisting}

Após iniciar, acesse a documentação interativa no Swagger em: \url{http://localhost:5000/swagger}

Para testar os endpoints com autenticação, clique em \textbf{Authorize} no Swagger e informe o token obtido no login no formato \texttt{Bearer \{token\}}.

% ==============================================================================
\section{Testes Automatizados e Integração Contínua}
% ==============================================================================

O projeto também possui o projeto \texttt{AcademicEvents.Tests}, criado com xUnit e Moq para validar as regras de negócio da camada Application sem depender do PostgreSQL. Essa abordagem permite testar os services diretamente, usando mocks das interfaces de repository definidas no Domain.

\begin{lstlisting}[language=bash]
dotnet test AcademicEvents.sln
\end{lstlisting}

\begin{longtable}{L{4cm} L{10.5cm}}
\toprule
\rowcolor{tablehead}
\color{white}\textbf{Service testado} & \color{white}\textbf{Cenários cobertos} \\
\midrule
\endhead
AuthService & Email duplicado no cadastro e login com credenciais inválidas \\
\rowcolor{tablealt}
EventService & Data final inválida, status inválido, organizador inválido e edição por outro usuário \\
RegistrationService & Evento inexistente e inscrição duplicada \\
\rowcolor{tablealt}
CommentService & Evento inexistente e remoção por usuário que não é autor \\
ReactionService & Evento inexistente e reação duplicada \\
\bottomrule
\end{longtable}

Além dos testes locais, foi criado o workflow \texttt{.github/workflows/ci.yml}. Ele executa restauração de dependências, compilação em modo Release e testes automatizados em pushes para \texttt{develop}, \texttt{main} e pull requests. Assim, cada alteração enviada ao repositório passa por uma verificação mínima antes de ser integrada.

% ==============================================================================
\section{Conclusão}
% ==============================================================================

A Academic Events API atende todos os requisitos obrigatórios do trabalho final: autenticação JWT, cinco entidades com relacionamentos, banco PostgreSQL via EF Core, arquitetura em cinco projetos separados, DTOs request/response, services com regras de negócio, repositories com repository pattern, Swagger/OpenAPI e versionamento via Git.

Além dos requisitos mínimos, o grupo implementou:

\begin{itemize}[noitemsep]
  \item Cinco exceções customizadas com semântica clara de negócio
  \item Filtragem de eventos por status e por organizador
  \item Validações em duas camadas (service + banco) para evitar dados inconsistentes
  \item Testes automatizados com xUnit e Moq para os services principais
  \item Workflow de CI com restore, build e testes automatizados
  \item Documentação XML em português em todos os arquivos
  \item Diagramas C4 nos quatro níveis (contexto, container, componente e código)
  \item Roteiro completo de testes manuais com cenários de sucesso e de erro
\end{itemize}

O projeto demonstra domínio prático dos conceitos de C\# e ASP.NET Core abordados no módulo, resultando em uma API funcional, bem organizada e pronta para ser consumida por qualquer aplicação cliente.

\end{document}
